\section{Preliminaries}
\label{sec:prelim}

\stitle{Semantic \textsf{ORDER BY}.}
Let $D=\{x_1,\dots,x_n\}$ be a collection of $n$ items, such as documents, passages, or records with unstructured text fields. 
Let $\phi$ be a natural-language instruction that specifies a desired ordering criterion, e.g., \emph{``order the reviews from most to least helpful for a purchase decision''}. 
A semantic \textsf{ORDER BY} operator returns a permutation $\Pi = [\pi_1, \cdots, \pi_n]$ of $D$ that is consistent with the order criterion induced by $\phi$. We assume the permutation $\Pi$ is in non-decreasing order throughout this paper. 
Unlike conventional \textsf{ORDER BY}, where the ordering key is a deterministic attribute value that can be compared directly by database engine, the semantic \textsf{ORDER BY},  defined through natural-language semantics, may depend on implicit notions such as relevance, helpfulness, specificity, sentiment, or evidential support, and therefore cannot be evaluated by a traditional comparator. We use a large language model (LLM) as an ordering oracle for this semantic relation. 

This LLM-based formulation introduces three challenges for physical algorithm design. First, invoking an LLM is expensive in both latency and token cost, with latency dominating conventional in-memory computation. Second, each invocation can inspect only a bounded number of items due to the effective context length limit of LLMs. Third, the intrinsic uncertainty of LLM reasoning can violate antisymmetry and transitivity due to positional bias. 
We design the merge operator to reduce input-token cost. We assess ranking quality separately because different call schedules can elicit different judgments. 



% This LLM-based semantic ORDER BY manifests three properties regarding algorithm design. First, an invocation is expensive:
% its token and latency cost dominate conventional
% query-processing costs, therefore the objective is to compute $\Pi$ with as
% few invocations as possible. Second, a single invocation can observe
% only a bounded number of items, so no call can rank $D$ at once and
% the global order must be assembled from partial rankings, which is
% precisely the access pattern of external merge sort.

\stitle{Listwise Ranking with LLMs.}
% Existing studies on LLM-based ranking propose three categorical methods with different LLM prompting. Pointwise methods~\cite{} evaluate each item independently and then sort the items by the resulting scores. The algorithm is easy to parallel, but the scores are difficult to calibrate across independent calls. 
% Pairwise methods~\cite{} request the model to compare two items at a time and then aggregate the resulting comparisons through sorting or graph-based algorithms. They can provide accurate relative judgments but require $O(n^2)$ invocations to construct a full ranking.
Listwise methods~\cite{DBLP:conf/sigir/ZhuangZKZ24, DBLP:conf/emnlp/0001YMWRCYR23, DBLP:journals/corr/abs-2603-17223, DBLP:journals/corr/abs-2305-02156} submit a batch of items in a single prompt and ask the model to return their identifiers in sorted order, thereby amortizing one invocation over many comparisons.
This paper adopts listwise prompting as the only LLM invocation primitive.
Formally, $\textsc{ListwiseRank}(B,\phi)$ takes a batch $B$ of $m$ items, where $m$ is determined by the LLM's effective context window and item serialization length, and returns an ordered sequence $\mathcal{O}$ to be a permutation of $B$ under instruction $\phi$. Throughout this paper, we assume $m\ge 4$, the minimum listwise batch capacity required by our loser-tree construction.
Each request carries the instruction and the serialized items in its batch. With the same template and similar batch widths, call count is a useful cost proxy. During a merge, batch widths and the number of times an item is resubmitted can differ across plans. Section~\ref{sec:cost} models both sources of input-token cost.

% \paragraph{Serialization and output repair}
% Figure~\ref{fig:prompt} shows the prompt template. Each item is
% serialized under a temporary identifier
% $[1],\dots,[|B|]$, and the model is instructed to output only the
% identifier sequence. The identifiers make the call auditable: an LLM
% may drop, duplicate, or fabricate items in its output, and the merger
% detects these defects by comparing the returned identifiers against
% the submitted ones. Identifiers missing from the response are
% appended in their input order, duplicates after the first occurrence
% are discarded, and identifiers that were never submitted are ignored,
% so every call returns a true permutation of its input batch.
% %
% Every listwise call submits at most $m$ items together with a fixed
% instruction, so calls carry a comparable token volume. The number of
% listwise calls is therefore a linear proxy for token cost and
% end-to-end latency, and all cost models in this paper
% (\cref{sec:sw-cost} and~\cref{sec:lt-cost}) minimize this count.

% ---------------------------------------------------------------------
% Prompt figure (requires \usepackage{listings})
% ---------------------------------------------------------------------
% \begin{figure}[t]
% \lstset{
%   basicstyle=\ttfamily\footnotesize,
%   frame=single,
%   breaklines=true,
%   breakindent=0pt,
%   columns=fullflexible,
%   keepspaces=true,
%   escapeinside={(*}{*)}
% }
% \begin{lstlisting}
% SYSTEM: You are an expert ranking assistant. Follow the instruction exactly and output only the requested format.

% USER: Rank the following (*$|B|$*) items by the criterion below. Criterion: (*$\phi$*)

% [1] (*item $x_1$*)
% [2] (*item $x_2$*)
% ...
% [(*$|B|$*)] (*item $x_{|B|}$*)

% Rank ALL (*$|B|$*) items from the one that best satisfies the criterion to the one that satisfies it least. Output only the identifiers in the format: [3] > [1] > ... > [7].
% Do not add, drop, or repeat any identifier. Do not output any explanation.
% \end{lstlisting}
% \caption{Prompt template of
% $\textsc{ListwiseRank}(B,\phi)$.}
% \label{fig:prompt}
% \end{figure}
